\section{서론}

한국지구과학회지(JKESS)는 지구과학 및 지구과학교육 관련 연구를 다루므로, 서론에서는 연구가 어떤 지구 시스템, 지구과학 개념, 관측 자료, 교육 문제와 연결되는지 명확히 제시하세요.

대기, 해양, 지질, 천문, 환경, 지구시스템 과학, 지구과학교육 등 원고의 영역을 초반에 분명히 하며, 선행연구 검토에서는 본 연구의 필요성과 차별성을 드러냅니다.

\subsection{학술지 및 투고 성격 확인}

한국지구과학회지는 한국지구과학회가 발행하는 학술지로, 지구과학 및 지구과학교육 관련 연구를 다룹니다. 투고 전 최신 투고규정에서 원고 종류, 초록과 주요어 길이, 그림과 표 작성 언어, 참고문헌 양식, 저자소개 요구사항을 확인하세요.

\subsection{연구 목적}

연구 목적 또는 연구 문제를 입력하세요. 지구과학 연구라면 자료와 분석 대상의 범위를, 지구과학교육 연구라면 학습자, 수업, 평가 또는 교육과정 맥락을 함께 밝혀 주세요.
